\section{A local rewrite description}

In this section we describe $\Fold_m$ by local rewriting on finitely supported
configurations.  This gives a concrete normalization procedure before we pass
to inverse limits and sequence-level dynamics.

\subsection{Value-preserving local rules}

Let
\[
\mathcal D
:=
\left\{
a=(a_1,a_2,\ldots):
a_k\in\NN,\ \supp(a)\text{ finite}
\right\},
\]
and define
\[
\Val(a):=\sum_{k\ge 1} a_kF_{k+1}.
\]
We write $e_k$ for the $k$th standard basis vector and
\[
\iota_m(\omega):=(\omega_1,\ldots,\omega_m,0,0,\ldots)\in\mathcal D
\]
for the natural embedding of a window into $\mathcal D$.

The rewrite system acts on $\mathcal D$ through the following local rules:
\begin{align}
e_k+e_{k+1}&\longrightarrow e_{k+2} && (k\ge 1), \label{eq:rewrite-adjacent}\\
2e_1&\longrightarrow e_2, \label{eq:rewrite-base1}\\
2e_2&\longrightarrow e_1+e_3, \label{eq:rewrite-base2}\\
2e_k&\longrightarrow e_{k-2}+e_{k+1} && (k\ge 3). \label{eq:rewrite-duplicate}
\end{align}
Each rule is local and preserves $\Val$ because it is exactly the Fibonacci
identity
\[
F_{k+1}+F_{k+2}=F_{k+3},
\qquad
2F_2=F_3,
\qquad
2F_3=F_2+F_4,
\qquad
2F_{k+1}=F_{k-1}+F_{k+2}.
\]

\begin{lemma}[Fibonacci rigidity of value-preserving local weights]
\label{lem:fold-local-weight-rigidity-fibonacci}
Suppose a positive weight sequence $(w_k)_{k\ge 1}$ has the property that each
rewrite rule \eqref{eq:rewrite-adjacent}--\eqref{eq:rewrite-duplicate}
preserves the weighted sum
\[
\Val_w(a):=\sum_{k\ge 1} a_kw_k
\qquad (a\in\mathcal D).
\]
Then
\[
w_{k+2}=w_{k+1}+w_k
\qquad (k\ge 1).
\]
In particular, under the normalization $w_1=F_2=1$ and $w_2=F_3=2$, one has
$w_k=F_{k+1}$ for all $k\ge 1$.
\end{lemma}

\begin{proof}
Preservation of \eqref{eq:rewrite-adjacent} gives
\[
w_k+w_{k+1}=w_{k+2}
\qquad (k\ge 1),
\]
which is exactly the Fibonacci recursion.  The remaining rules are then
automatic consequences of that recursion:
\[
2w_1=w_2,\qquad 2w_2=w_1+w_3,\qquad 2w_k=w_{k-2}+w_{k+1}\ \ (k\ge 3).
\]
With the stated normalization, the recursion has the unique solution
$w_k=F_{k+1}$.
\end{proof}

\begin{lemma}[Local weight rigidity]
\label{lem:local-weight-rigidity}
Let
\[
A(a):=\sum_{k\ge 1} a_k,
\qquad
B(a):=\sum_{k\ge 1} ka_k,
\qquad
C(a):=a_2,
\]
and order triples lexicographically.  Then every rewrite step
$a\to a'$ generated by \eqref{eq:rewrite-adjacent}--\eqref{eq:rewrite-duplicate}
strictly decreases
\[
\mu(a):=(A(a),B(a),C(a)).
\]
In particular the rewrite system is strongly terminating.
\end{lemma}

\begin{proof}
Rules \eqref{eq:rewrite-adjacent} and \eqref{eq:rewrite-base1} reduce the total
digit count $A$ by one.  Rule \eqref{eq:rewrite-duplicate} for $k\ge 3$
preserves $A$ but changes $B$ by
\[
2k-\bigl((k-2)+(k+1)\bigr)=1,
\]
so $B$ decreases by one.  Rule \eqref{eq:rewrite-base2} preserves both $A$ and
$B$ but decreases $C=a_2$ by two.  Hence every rewrite step strictly decreases
$\mu(a)$.
\end{proof}

\subsection{Canonical normal forms}

\begin{proposition}[Termination, confluence, and normal forms]
\label{prop:rewrite-newman}
The rewrite system on $\mathcal D$ generated by
\eqref{eq:rewrite-adjacent}--\eqref{eq:rewrite-duplicate} is strongly
terminating and confluent.  Every $a\in\mathcal D$ therefore has a unique
normal form
\[
\NF(a)\in\Xfin
\]
satisfying
\[
\Val(\NF(a))=\Val(a).
\]
\end{proposition}

\begin{proof}
Strong termination is Lemma~\ref{lem:local-weight-rigidity}.  A configuration
is irreducible exactly when every digit is in $\{0,1\}$ and no adjacent pair
$11$ occurs; this is precisely the definition of $\Xfin$.

Fix $a\in\mathcal D$.  By termination, every maximal reduction sequence from
$a$ ends in some irreducible configuration $x\in\Xfin$.  Since each rewrite
preserves $\Val$, every such terminal configuration satisfies
\[
\Val(x)=\Val(a).
\]
If $x$ and $y$ are two terminal reductions of $a$, then $x,y\in\Xfin$ are
irreducible configurations with the same Fibonacci value.  Zeckendorf
uniqueness therefore forces $x=y$.  Thus every $a$ has a unique irreducible
descendant.  For a terminating rewrite system, uniqueness of irreducible
descendants is equivalent to confluence, so the system is confluent and the
normal form $\NF(a)$ is well defined.
\end{proof}

\begin{corollary}[Order-independent implementation of the fold]
\label{cor:fold-order-independent}
For every $\omega\in\Omega_m$,
\[
\Fold_m(\omega)=\pi_m(\NF(\iota_m(\omega))),
\]
and the right-hand side is independent of the order in which rewrites are
applied.
\end{corollary}

\begin{proof}
By Proposition~\ref{prop:rewrite-newman}, the normal form of $\iota_m(\omega)$
is the unique admissible Fibonacci digit string with value $N(\omega)$, namely
$Z(N(\omega))$.  Taking the first $m$ digits gives the claim.
\end{proof}

\begin{remark}[Two illustrative normalizations]
\label{rem:two-fold-examples}
The following two examples illustrate the order-independence of the rewriting
and the appearance of the overflow bit.

For $\omega=1110\in\Omega_4$ one may reduce
\[
\iota_4(\omega)=e_1+e_2+e_3\to 2e_3\to e_1+e_4,
\]
so $\NF(\iota_4(\omega))=\iota_4(1001)$ and
\[
\Fold_4(1110)=1001.
\]
If instead $\omega=1101\in\Omega_4$, then
\[
\iota_4(\omega)=e_1+e_2+e_4\to e_3+e_4\to e_5,
\]
hence
\[
\Fold_4(1101)=0000,
\qquad
N(1101)=F_6=F_{4+2}.
\]
This is the smallest example in which the visible word is zero while the whole
value survives as the overflow bit from
Corollary~\ref{cor:visible-value-overflow}.
\end{remark}

\begin{proposition}[Basic properties of $\Fold_m$]
\label{prop:fold-basic}
For every $m\ge 1$, the map $\Fold_m:\Omega_m\to X_m$ is well defined and
satisfies:
\begin{enumerate}[label=\textup{(\roman*)}]
\item $\Fold_m(\omega)\in X_m$ for every $\omega\in\Omega_m$;
\item $\Fold_m(x)=x$ for every $x\in X_m$;
\item $\Fold_m:\Omega_m\twoheadrightarrow X_m$ is surjective.
\end{enumerate}
\end{proposition}

\begin{proof}
Property (i) is immediate from Definition~\ref{def:fold-word}.  For (ii), if
$x\in X_m$ then $x$ is already an admissible Zeckendorf prefix, so
$Z(N(x))=(x_1,\ldots,x_m,0,0,\ldots)$ and hence $\Fold_m(x)=x$.  Surjectivity
follows from (ii): every $x\in X_m$ is a fixed point of $\Fold_m$.
\end{proof}

